\documentclass[12pt]{article}
\usepackage{amsmath, amssymb}
\usepackage{geometry}
\geometry{margin=1in}

\begin{document}

\section*{CSE400 – Fundamentals of Probability in Computing}

\subsection*{Lecture 8: Gaussian, Uniform, Exponential, and Gamma Random Variables}

\textbf{Instructor:} Dhaval Patel, PhD \\
\textbf{Date:} January 29, 2026

\hrule
\vspace{1em}

\section*{1. Lecture Outline (As Shown on Slides)}

The lecture is organized into the following sections:

\begin{enumerate}
    \item \textbf{Gaussian Random Variable}
    \begin{itemize}
        \item Definition and Properties
        \item Standard Forms: Q-function and Error Function
        \item Relation between $\Phi$-function and Q-function
        \item Example
    \end{itemize}

    \item \textbf{Types of Continuous Random Variable}
    \begin{itemize}
        \item Uniform Random Variable: Example
        \item Exponential Random Variable: Example
        \item Laplace Random Variable: Example
        \item Gamma Random Variable
        \begin{itemize}
            \item Graph and Special Cases
            \item Example
        \end{itemize}
    \end{itemize}

    \item Homework Problem
    \item Problem Solving
    \item In-Class Activity: Gaussian Density Estimation
\end{enumerate}

\begin{quote}
Only the topics that are \textbf{explicitly developed on the slides} are expanded below. \\
Topics listed in the outline but \textbf{not developed further on the slides} are not expanded.
\end{quote}

\hrule
\vspace{1em}

\section*{2. Expected Value and Moments (Material Appearing Before Gaussian RV)}

\subsection*{2.1 Expected Value of a Function of a Random Variable}

\subsubsection*{Continuous Random Variable (CRV)}

The expected value of a function $g(X)$ is given by:
\[
\mathbb{E}[g(X)] = \int g(x)\, f_X(x)\, dx
\]

\subsubsection*{Discrete Random Variable (DRV)}

The expected value of a function $g(X)$ is given by:
\[
\mathbb{E}[g(X)] = \sum_k g(x_k)\, P_X(x_k)
\]

These formulas are written explicitly on the slide titled \\
\textbf{“Expected value of g of RV: X”}

\subsection*{2.2 Linearity Properties of Expectation}

For constants $a$ and $b$:
\[
\mathbb{E}[aX + b] = a\,\mathbb{E}[X] + b
\]

For a finite sum of functions:
\[
\mathbb{E}\!\left[\sum_{k=1}^{N} g_k(X)\right]
=
\sum_{k=1}^{N} \mathbb{E}[g_k(X)]
\]

These properties are stated directly on the slide, with no proof or additional explanation.

\hrule
\vspace{1em}

\section*{3. Higher-Order Central Moments}

\subsection*{3.1 Third Central Moment – Skewness}

The third central moment is defined as:
\[
m_3 = C_s =
\frac{\mathbb{E}[(X - \mu_X)^3]}{\sigma_X^3}
\]

\textbf{Explanation exactly as written on the slide:}
\begin{itemize}
    \item It is a \textbf{measure of symmetry of the PDF about the mean ($\mu_X$)}
    \item Positive value $\rightarrow$ \textbf{Right}
    \item Negative value $\rightarrow$ \textbf{Left}
\end{itemize}

No additional explanation is given.

\subsection*{3.2 Fourth Central Moment – Kurtosis}

The fourth central moment is defined as:
\[
m_4 = C_k =
\frac{\mathbb{E}[(X - \mu_X)^4]}{\sigma_X^4}
\]

\textbf{Explanation exactly as written on the slide:}
\begin{itemize}
    \item Large value $\rightarrow$ the random variable will have a \textbf{large peak near the mean}
\end{itemize}

A sketch of a sharp peak is shown. No algebraic derivation is provided.

\hrule
\vspace{1em}

\section*{4. Gaussian Random Variable}

\subsection*{4.1 Definition}

A \textbf{Gaussian random variable} is defined as a random variable whose \textbf{PDF} can be written in the general form:
\[
f_X(x) =
\frac{1}{\sqrt{2\pi\sigma^2}}
\exp\!\left(
-\frac{(x - m)^2}{2\sigma^2}
\right)
\]

This definition appears under \\
\textbf{“Gaussian Random Variable – Definition and Properties”}

\subsection*{4.2 PDF and CDF Representation}

The slide shows:
\begin{itemize}
    \item The \textbf{PDF} of a Gaussian random variable
    \item The corresponding \textbf{CDF}
    \item A labeled example with $m = 3$ and $\sigma = 2$
\end{itemize}

The visual relationship shown:
\begin{itemize}
    \item Area under the PDF corresponds to the CDF
    \item Left-tail area corresponds to $\Phi(x)$
    \item Right-tail area corresponds to $Q(x)$
\end{itemize}

No algebraic explanation beyond the figure is included.

\hrule
\vspace{1em}

\section*{5. Standard Forms for Gaussian Random Variables}

\subsection*{5.1 Error Function}

The \textbf{error function} is defined as:
\[
\mathrm{erf}(x) =
\frac{2}{\sqrt{\pi}}
\int_0^x e^{-t^2}\, dt
\]

\subsection*{5.2 Complementary Error Function}

The \textbf{complementary error function} is defined as:
\[
\mathrm{erfc}(x)
= 1 - \mathrm{erf}(x)
=
\frac{2}{\sqrt{\pi}}
\int_x^{\infty} e^{-t^2}\, dt
\]

Both definitions are written explicitly on the slides.

\subsection*{5.3 $\Phi$-Function (CDF of Standard Normal)}

The $\Phi$-function is defined as:
\[
\Phi(x) =
\frac{1}{\sqrt{2\pi}}
\int_{-\infty}^{x}
\exp\!\left(-\frac{t^2}{2}\right) dt
\]

It is explicitly labeled as the \textbf{CDF of the Standard Normal}.

\subsection*{5.4 Q-Function (Gaussian Tail Function)}

The Q-function is defined as:
\[
Q(x) =
\frac{1}{\sqrt{2\pi}}
\int_x^{\infty}
\exp\!\left(-\frac{t^2}{2}\right) dt
\]

It is explicitly labeled as the \textbf{Gaussian Tail Function}.

\subsection*{5.5 Relationship Between $\Phi$ and Q}

A note on the slide states:
\[
Q(x) = 1 - \Phi(x)
\]

No derivation is provided.

\hrule
\vspace{1em}

\section*{6. Relation Between $\Phi$-Function and Q-Function for a Gaussian RV}

\subsection*{6.1 Evaluating the CDF}

For a Gaussian random variable $X \sim \mathcal{N}(m,\sigma^2)$:
\[
F_X(x) =
\int_{-\infty}^{x}
\frac{1}{\sqrt{2\pi\sigma^2}}
\exp\!\left(
-\frac{(y - m)^2}{2\sigma^2}
\right) dy
\]

Change of variable shown on the slide:
\[
\int_{-\infty}^{\frac{x-m}{\sigma}}
\frac{1}{\sqrt{2\pi}}
\exp\!\left(-\frac{t^2}{2}\right) dt
\]

Therefore:
\[
F_X(x) =
\Phi\!\left(\frac{x-m}{\sigma}\right)
\]

All steps appear explicitly on the slide.

\subsection*{6.2 Evaluating Tail Probabilities}

\[
\Pr(X > x) =
\int_{\frac{x-m}{\sigma}}^{\infty}
\frac{1}{\sqrt{2\pi}}
\exp\!\left(-\frac{t^2}{2}\right) dt
\]

\[
= Q\!\left(\frac{x-m}{\sigma}\right)
\]

\subsection*{6.3 Key Identities}

\[
Q(x) = 1 - \Phi(x)
\]

\[
F_X(x) = 1 - Q\!\left(\frac{x-m}{\sigma}\right)
\]

\subsection*{6.4 Notes Provided on Slides}

\begin{itemize}
    \item To evaluate the \textbf{CDF} of a Gaussian random variable, evaluate the \textbf{$\Phi$-function} at the points.
    \item The \textbf{Q-function} is more natural for evaluating probabilities of the form $\Pr(X > x)$.
\end{itemize}

\hrule
\vspace{1em}

\section*{7. Gaussian Random Variable – Worked Example}

\subsection*{7.1 Given PDF}

\[
f_X(x) =
\frac{1}{\sqrt{8\pi}}
\exp\!\left(-\frac{(x+3)^2}{8}\right)
\]

\subsection*{7.2 Identification of Parameters}

Comparing with the Gaussian PDF:
\[
f_X(x) =
\frac{1}{\sqrt{2\pi\sigma^2}}
\exp\!\left(-\frac{(x-m)^2}{2\sigma^2}\right)
\]

We identify:
\[
m = -3,\quad \sigma^2 = 4,\quad \sigma = 2
\]

\subsection*{7.3 Gaussian Relations Used}

\[
F_X(x) =
\Phi\!\left(\frac{x-m}{\sigma}\right)
\]

\[
\Pr(X > x) =
Q\!\left(\frac{x-m}{\sigma}\right)
\]

\subsection*{7.4 Probability Calculations}

\paragraph{(1) $\Pr(X \le 0)$}

\[
\Pr(X \le 0) =
F_X(0)
\]

\[
= \Phi\!\left(\frac{0 - (-3)}{2}\right)
= \Phi\!\left(\frac{3}{2}\right)
\]

\[
= 1 - Q\!\left(\frac{3}{2}\right)
\]

\paragraph{(2) $\Pr(X > 4)$}

\[
\Pr(X > 4) =
Q\!\left(\frac{4 - (-3)}{2}\right)
= Q\!\left(\frac{7}{2}\right)
\]

\end{document}
